\subsection{The exact sign schedule}

The complete raw body has 359 coordinates and 307 quadratic rows. The final guarded source adds 36 scalar nonnegativity predicates and twelve product predicates, each represented with three new integer witnesses. It therefore has 503 coordinates and 355 rows. The later exact monic projection deletes 125 coordinates and 125 defining rows, leaving the measured 378-variable, 230-residual equation.

The sign implications are ordered as follows.
\begin{enumerate}
\item The 36 scalar predicates and literal polynomial implications establish nonnegativity of every semantic frame atom, positivity of all fifteen requested indices, and the complete natural domain of the private anchor. The aggregate quotient $Q$ is among these retained conditions. All six proposed exact-power outputs have independently established signs at this stage. Four of these output predicates are explicitly retained to prevent dependence on remainder pairs that will only be interpreted later.
\item The full oriented-pair aggregate theorem now identifies both coordinates of all fifteen seed pairs. It supplies thirty signs and their exact mathematical meanings simultaneously.
\item Decode the exact powers in their dependency order: first the requests with bases $F^2+1$ and $B\ge2$, then the request with the now established base $N$, and then the three denominator requests at base $q$. For each base $1\le n<W$, $T_n=2An-n^2-1>0$ follows from the frame. The initially unrestricted quotient in its power row becomes nonnegative by \eqref{eq:recpower}. No sign of a not-yet-interpreted binomial is used.
\item Restore the twelve remainder pairs. Their actual rows have $r+s=d-1$, where $d=B\ge2$ for the factorial, or $d$ is one of the now certified positive powers of $q$. The predicate $rs\ge0$ then forces $r,s\ge0$: if one were negative and their product nonnegative, their sum would be negative. A separate anchor pair is unnecessary because $d\ge2$ is already established. These twelve product conditions replace twenty-four scalar sign requirements.
\item The single exact-binomial output slack is $q-B-1\ge0$, since $q=B^{K^2}$, $B\ge2$, and $K\ge2$. The common-root quotient is nonnegative because its numerator is $2A(q+1)-(q+1)^2-1>0$ and its modulus is positive. Each of the eleven modular digit quotients is nonnegative because its positive canonical Pell residue is congruent to the indicated least nonnegative residue. This argument precedes its later factorial or range interpretation.
\item The checked polynomial implication ledger, now using these recovered signs, restores all 359 original body coordinates to their natural domains. The whole original arithmetic-source theorem is applied only after that restoration.
\end{enumerate}

The paired remainder predicates are bound to twelve specific rows: the factorial pair, the exact binomial pair, eight local-column range pairs, the index range pair, and the derived-label range pair. Their actual denominators, coordinate maps, and guard equations appear in the signed certificate. There are 22 initially deferred quotient/modulus coordinates: six exact-power quotients, four distinct base moduli, eleven modular digit quotients, and the common-root quotient. The separately delayed output slack is recorded as well. These are concrete finite lists, not an informal permission to omit arbitrary signs.

\subsection{Why the unrestricted modular quotients cause no extra branch}

Let $H_e=\chi_A(e)+(q+1-A)\psi_A(e)$. Once the principal pair is canonical, $H_e>0$, since $\chi_A(e)>(A-1)\psi_A(e)$ and $q\ge1$. In the exact digit row the chosen residue $r+dB$ belongs to $[0,qd)$; in a range row $r\in[0,d)\subset[0,Nd)$. An integral quotient expressing a positive integer as that least nonnegative residue plus a multiple of the positive modulus cannot be negative. This restores the ordinary natural digit graph without presupposing the digit's intended binomial value.
